\chapter{Stratégie marketing}
\label{ch:marketing}

\section{Introduction}

Même la place de marché la plus fiable échoue sans notoriété, adoption et rétention. Ce chapitre définit la stratégie marketing de MABRID : positionnement de marque, acquisition client, canaux numériques, croissance communautaire et programmes de fidélité---alignés sur la vision commerciale centrée sur la confiance et sans recours à des tactiques de croissance trompeuses.

\section{Image de marque}

\subsection{Identité de marque}

La marque MABRID exprime la \emph{découverte locale de confiance} : accessible, professionnelle, ancrée au Maroc et moderne. L'identité visuelle (à fournir sous forme de figures) utilise des couleurs, une typographie et des combinaisons logo cohérentes sur tous les points de contact.

\figureplaceholder{H}{figures/ch07-brand-guidelines.pdf}{0.75}{Aperçu des directives de marque MABRID (insérer l'image du brand board)}

\subsection{Promesse de marque}

« Trouvez des entreprises locales auxquelles vous pouvez faire confiance. » La promesse est opérationnalisée par la vérification, la modération et des politiques transparentes---le marketing ne doit pas surestimer les capacités de la plateforme.

\subsection{Architecture de marque}

Marque maîtresse MABRID avec sous-marques pour les communautés par ville (p. ex. MABRID Casablanca) et les niveaux premium (MABRID Pro pour les vendeurs).

\section{Positionnement}

\subsection{Énoncé de positionnement}

Pour les consommateurs et PME marocains qui peinent avec une découverte locale fragmentée, MABRID est la place de marché de confiance qui connecte des entreprises vérifiées aux acheteurs par une recherche pertinente par ville, des avis et une communauté---contrairement aux petites annonces généralistes ou à la vente sociale informelle, car MABRID investit dans la gouvernance, la sécurité et la conformité.

\subsection{Différenciation concurrentielle}

Le tableau~\ref{tab:positioning} contraste les dimensions de positionnement.

\begin{table}[H]
  \centering
  \caption{Dimensions de positionnement concurrentiel}
  \label{tab:positioning}
  \begin{tabularx}{\textwidth}{@{}lXX@{}}
    \toprule
    \textbf{Dimension} & \textbf{Petites annonces généralistes} & \textbf{MABRID} \\
    \midrule
    Signaux de confiance & Limités & Vérification, avis, modération \\
    Pertinence locale & Variable & Communautés par ville et cartes \\
    Outils vendeurs & Annonces de base & Vitrine, analyses, premium \\
    Posture de conformité & Souvent minimale & Suite juridique publiée, alignement CNDP \\
    \bottomrule
  \end{tabularx}
\end{table}

\section{Acquisition client}

\subsection{Acquisition acheteurs}

Les canaux comprennent l'optimisation pour les moteurs de recherche sur les requêtes locales, l'optimisation des boutiques d'applications, les partenariats avec des influenceurs lifestyle et food, et les relations publiques autour des initiatives de confiance (jalons de vérification, lancements par ville).

\subsection{Acquisition vendeurs}

Prospection terrain auprès des associations professionnelles, événements des chambres de commerce, ateliers de compétences numériques et incitations de parrainage pour les comptables et agences web servant les PME.

\subsection{Entonnoir d'acquisition}

Notoriété $\rightarrow$ considération (pages d'atterrissage, études de cas) $\rightarrow$ inscription $\rightarrow$ activation (première annonce ou première recherche) $\rightarrow$ rétention. Les métriques d'entonnoir sont revues hebdomadairement pendant les phases de croissance.

\figureplaceholder{H}{figures/ch07-marketing-funnel.pdf}{0.85}{Entonnoir d'acquisition client pour acheteurs et vendeurs}

\section{Marketing numérique}

\subsection{Recherche et contenu}

Le marketing de contenu localisé (blog, guides : « Meilleurs plombiers à Rabat ») génère du trafic organique tout en apportant une valeur réelle aux utilisateurs. Le contenu respecte les normes publicitaires et de protection des consommateurs.

\subsection{Médias payants}

Les campagnes sociales et de recherche payantes ciblent des mots-clés par ville avec des budgets plafonnés et des tests créatifs. Les pages d'atterrissage reflètent les promesses publicitaires pour éviter le rebond et les plaintes réglementaires.

\subsection{Courriel et cycle de vie}

Le courriel transactionnel (vérification, alertes) est distingué du courriel marketing (opt-in uniquement). Les campagnes de cycle de vie réengagent les acheteurs inactifs et éduquent les nouveaux vendeurs sur les avantages de la vérification.

\section{Stratégie réseaux sociaux}

Plateformes prioritaires : Instagram, Facebook, TikTok, LinkedIn (vendeurs B2B) et WhatsApp Business pour le support vendeurs (pas de spam acheteur non sollicité). Les community managers répondent en arabe et en français. Les playbooks de communication de crise s'intègrent à la confiance et sécurité.

\section{Programme de parrainage}

Les incitations de parrainage bilatérales récompensent acheteurs et vendeurs invitant des participants vérifiés. Les contrôles anti-fraude empêchent la création artificielle de comptes. Les conditions de parrainage sont publiées pour éviter la perception de systèmes pyramidaux.

\section{Rétention client}

\subsection{Rétention acheteurs}

Recherches enregistrées, favoris, recommandations personnalisées (avec contrôles de confidentialité), engagement communautaire par ville et notifications push pour les entreprises suivies augmentent les visites de retour.

\subsection{Rétention vendeurs}

Tableaux de bord de métriques de succès, support réactif, études de cas ROI premium et prestige du badge de vérification réduisent l'attrition. Les entretiens de sortie informent la feuille de route produit.

\subsection{Principes de fidélité}

La fidélité se gagne par la fiabilité, pas par des dark patterns. Le désabonnement et la suppression de compte restent accessibles selon la politique de confidentialité.

\section{Intégration avec la stratégie de croissance}

Le marketing s'aligne sur les phases du chapitre~\ref{ch:roadmap} : investissement local important dans les villes d'ancrage avant les dépenses médias nationales. Le budget marketing en pourcentage des revenus se resserre à mesure que la part organique croît.

\section{KPI marketing}

\begin{table}[H]
  \centering
  \caption{Indicateurs de performance marketing}
  \label{tab:marketing-kpis}
  \begin{tabularx}{\textwidth}{@{}lX@{}}
    \toprule
    \textbf{KPI} & \textbf{Objectif} \\
    \midrule
    CAC par canal & Allouer les dépenses efficacement \\
    Taux de conversion (visite $\rightarrow$ inscription) & Optimiser les expériences d'atterrissage \\
    Taux d'activation vendeurs & Mesurer la qualité d'intégration \\
    Rétention acheteurs à 30 jours & Évaluer l'adéquation produit-marché \\
    Part de voix (réseaux sociaux) & Suivre la notoriété de marque \\
    Net Promoter Score & Satisfaction globale \\
    \bottomrule
  \end{tabularx}
\end{table}

\section{Conclusion du chapitre}

La stratégie marketing de MABRID renforce---jamais ne compromet---son positionnement de confiance. Une image de marque disciplinée, une acquisition localisée, des campagnes numériques éthiques et un marketing produit axé sur la rétention construisent ensemble une base d'utilisateurs durable capable de soutenir l'expansion nationale et la viabilité financière à long terme.
